%------------------------------------------------------------------------------
% Voorbeeld presentatie HoGent - huisstijl 2018
%------------------------------------------------------------------------------
\documentclass[aspectratio=169]{beamer}

% Je kan het lettertype iets vergroten door hierboven optie ``14pt'' toe te
% voegen.

%==============================================================================
% Aanloop
%==============================================================================

%---------- Vormgeving --------------------------------------------------------

\usetheme{hogent}

% Kies hieronder een achtergrondkleur
\usecolortheme{hgwhite} % witte achtergrond, zwarte tekst
%\usecolortheme{hgblack} % zwarte achtergrond, witte tekst

%---------- Packages ----------------------------------------------------------

\usepackage[dutch]{babel}      % Nederlandse taal: splitsingen, enz.

\usepackage{booktabs}          % Mooie tabellen
\usepackage{multirow,multicol} % Tabelcellen samenvoegen
\usepackage{eurosym}           % Euro symbool

\usepackage{lipsum}             % Vultekst (enkel nuttig voor dit voorbeeld)

%---------- Commando-definities -----------------------------------------------

%---------- Info over de presentatie ------------------------------------------

\title[Korte titel]{Titel van de presentatie.}
\subtitle{subtitel}
\author[BVV]{Bert {Van Vreckem} (\href{mailto:bert.vanvreckem@hogent.be}{bert.vanvreckem@hogent.be})}
\date{\today}

%==============================================================================
% Inhoud presentatie
%==============================================================================

\begin{document}

{
\setbeamertemplate{background}[imgletter]{img/blurred-background-cellphone-coffee-842554.jpg}{H}
% Bron afbeelding: https://tex.stackexchange.com/questions/63146/negative-font-overlayed-with-a-photo

\begin{frame}
  \maketitle
\end{frame}
}

\begin{frame}
  \frametitle{Inhoud.}

  \tableofcontents
\end{frame}

\section{Eerste deel.}

\begin{frame}
  \frametitle{Inleiding.}
  
  \lipsum*[2]
\end{frame}

{
\setbeamertemplate{background}[imgletter]{img/blurred-background-cellphone-coffee-842554.jpg}{W}

\begin{frame}

{\huge \textbf{Voorbeeld van tekst door een afbeelding}}

\end{frame}
}

\subsection{Subsectie 1.1.}

\begin{frame}
  \frametitle{Titel.}

  \begin{itemize}
  \item Lijn één
  \item Lijn twee
  \item Lijn drie
  \end{itemize}
\end{frame}

\subsection{Subsectie 1.2.}

\begin{frame}
  \frametitle{Twee kolommen.}

  \begin{columns}[c]

  \column{.5\textwidth}
    \begin{itemize}
    \item Lijn 1
    \item Lijn 2
    \item Lijn 3
    \end{itemize}

  \column{.5\textwidth}
    \begin{itemize}
    \item Lijn 1
    \item Lijn 2
    \item Lijn 3
    \end{itemize}

  \end{columns}
\end{frame}

\begin{frame}
  \frametitle{Kleuren}
  
  \textbf{\textcolor{hgdarkgreen}{tekst}
  \textcolor{hgpink}{tekst}
  \textcolor{hgochre}{tekst}
  \textcolor{hgorange}{tekst}
  \textcolor{hgpurple}{tekst}
  \textcolor{hgblue}{tekst}
  \textcolor{hglightgreen}{tekst}
  \textcolor{hgbrown}{tekst}
  \textcolor{hggrey}{tekst}
  \textcolor{hgyellow}{tekst}}

  \textbf{\colorbox{hgdarkgreen}{tekst}
  \colorbox{hgpink}{tekst}
  \colorbox{hgochre}{tekst}
  \colorbox{hgorange}{tekst}
  \colorbox{hgpurple}{tekst}
  \colorbox{hgblue}{tekst}
  \colorbox{hglightgreen}{tekst}
  \colorbox{hgbrown}{tekst}
  \colorbox{hggrey}{tekst}
  \colorbox{hgyellow}{tekst}}

\end{frame}

\section{Tweede deel.}

\begin{frame}
  \frametitle{Het tweede deel.}
  
  \begin{table}
    \begin{tabular}{lll}
      \toprule
                     & \textbf{lorem ipsum} & \textbf{lorem ipsum} \\
      \midrule
      dolor sit amet & $3.255$              & $3.255$ \\
      \midrule
      lorem ipsum    & $425,43 (13,07\%)$   & $425,43 (13,07\%)$ \\
      \midrule
      dolor sit amet & $0,00$               & $1.273,31$ ($45\%$ -- gemiddeld) \\
      \midrule
      \textbf{{\small Verhouding kost}} & \textbf{$1,53$} & \textbf{$2,72$}\\
      \bottomrule
    \end{tabular}
    
    \label{tab:voorbeeld}
    \caption{Deze tabel heeft een bijschrift}
  \end{table}
\end{frame}

\subsection{Sectie 2.1.}

\begin{frame}
  \frametitle{Afbeeldingen.}
  \begin{figure}
    \caption{Deze afbeelding heeft een bijschrift.}
    \includegraphics[height=.8\textheight]{img/beverage-coffee-computer-984536.jpg}
    % Bron afbeelding: https://www.pexels.com/photo/hand-on-cup-of-coffee-984536/
    \label{img:voorbeeld}
  \end{figure}
\end{frame}

\begin{frame}
  \includegraphics[width=\textwidth]{img/beverage-coffee-computer-984536.jpg}
\end{frame}


{
\setbeamertemplate{background}[imgletter]%
  {img/beverage-coffee-computer-984536.jpg}{T}

\begin{frame}

  {\huge \textbf{This is a famous quote that says something about my point I want to make.}}
  
  \bigskip
  
  \textbf{Name Quotemaker}

\end{frame}
}

\begin{frame}[fragile]
  \texttt{monogespatieerd lettertype met ligaturen -> <= >=}

  \begin{verbatim}
  a <- 1
  b <- 2
  a + b
  \end{verbatim}
\end{frame}

\end{document}
